\clearpage\subsection{Prompt assembly for the boundary experiment}
All four prompts share the following literal strings. LF below denotes a single newline character. Soft line wrapping in the paper is typographical.
\paragraph{Common lead.}\begin{quote}\small
Use three internal checkpoints in this one response.
\end{quote}

\paragraph{Operations in fixed order.}
\paragraph*{Operation 1.}\begin{quote}\small
Analyze the requirements and edge cases and form a concise plan.
\end{quote}

\paragraph*{Operation 2.}\begin{quote}\small
Implement the requested function using that plan.
\end{quote}

\paragraph*{Operation 3.}\begin{quote}\small
Check the implementation against the requirements and revise it if needed.
\end{quote}

\paragraph{Common ending.}\begin{quote}\small
Do not output checkpoint notes or commentary. Return exactly one fenced Python program.
\end{quote}

\texttt{LABEL} denotes a stage heading; \texttt{OPERATION} denotes its exact operation sentence above. Role-labelled conditions use PLANNER, IMPLEMENTER, REVIEWER; neutral conditions use STAGE 1, STAGE 2, STAGE 3.
\begin{enumerate}
\item \textbf{Format the stages.} Heading-style stages use \texttt{[LABEL] OPERATION}, with LF between stages. Prose-style stages use \texttt{LABEL: OPERATION}, with one space between stages.
\item \textbf{Assemble the instruction block.} Join the common lead, the three formatted stages and the common ending. Put one LF before and after the stages.
\item \textbf{Assemble the request.} Concatenate \texttt{TASK}, LF, the original task instruction, two LF characters, the instruction block, LF, \texttt{FULL SUPPLIED CONTEXT}, LF and the complete prepared context, in that order.
\end{enumerate}
The common execution instruction is the same as above. Only stage labels and separators differ among conditions.
